\cotefile{15}
\scanpage{195}
2°) $X_{\deleted{\unsure{\mathrm{réd}}}}$ domine $\operatorname{Spec}\mathbf Z$, i.e. $X$ contient des pts de car. résiduelle nulle.

Alors \deleted{\unsure{en plus} \inserted{\unsure{dispose}} \deleted{\unsure{de}} \unsure{de ceci} (et éventuellement \unsure{cristalline} \unsure{le cas d'inégales car.}) on a, pour \unsure{tt pt}}

\deleted{\unsure{ét. (crist.)}}

(i) $\xi_\ell\in\unsure{\mathbf T}(\mathbf Z_\ell)$ \uline{\unsure{tt}} $\ell\in\mathbf P$, $\xi_\ell$ provenant d'un \unsure{élém.} de $T_{\unsure{\alpha}}(\widetilde{\mathbf Z}_{\unsure{\ell\mathbf Z}})$.

$\xi_\ell$ dépend du choix d'un pt géom. quelconque de $X$, et est indép. de ce choix modulo opération du $\pi_1$, notation $\xi_{\bar x,\ell}:=\xi_{X,\ell}$.

(ii) Si $X$ a des points $x$ de car. $p>0$, alors on dispose également des $\xi_{x,p}$, $\xi_{x,p}^{0}$.

\unsure{Connexions} : \unsure{relation}, \deleted{\unsure{de la}} si $y$ est un pt de car. $p>0$ et \deleted{\unsure{$x$}} un pt de car. $0$, entre $\xi_{y,p}$ et $\xi_{x,p}\otimes_{\mathbf Z_p}W(k)$ [$y$ pt \unsure{à val.} sur un corps parfait $k$].

\uline{Ils ne sont en général pas semi-simples}, sans doute. \deleted{\unsure{Ils}} \unsure{Si} \deleted{\unsure{sont semi-simples}} leurs \unsure{poids}, ils sont cependant semi-simples \unsure{par le gr.}, \deleted{\unsure{et les $\xi_p$ doivent}} \deleted{\illegible{6}}.
\margin{\uuline{\deleted{\unsure{Vérifier}}}$\longrightarrow$}

(iii) Si $x$ est un pt \unsure{à val. ds.} un corps \inserted{$k$} de \uline{car. $0$}, on a $\xi_{x,\infty}\in T_{\unsure{\alpha}}(k)$ par Hodge-de Rham. Il n'y a pas lieu de \inserted{Avec ses \textcircled{\uuline{2}} filtrations \deleted{\unsure{algébriques}} !}

\margin{L'une est celle \unsure{qui} donne comme gr. associé les $\unsure{H^i}(x,\unsure{\Omega^j})$, l'autre (arithmétique) \deleted{\unsure{par}} \unsure{provient de} la filtration canonique des motifs. (Elle est \uline{triviale} pour les motifs simples) La première filtration est \unsure{originelle} au foncteur Hodge-De Rham, bien faire attention qu'elle ne provient pas d'une filtration des motifs, c'est une filtration du \uline{\unsure{foncteur}} !}
\margin{194}
